% =====================================================================
\chapter{Casos de estudio: moldeo por inyección y hormigón}
\label{tema:5}
% =====================================================================

Este tema reúne las piezas de los anteriores en dos casos completos: el ajuste de la inyectora, que se ha ido desarrollando desde el \cref{tema:1}, y la formulación de un hormigón con menos emisiones de CO$_2$. Ambos usan datos sintéticos generados por los simuladores de la Parte~5 de la práctica; la revisión de Fernandes \emph{et al.}~\cite{fernandes} sobre modelado y optimización del moldeo por inyección y el trabajo de Ament \emph{et al.}~\cite{ament} sobre hormigón sostenible con optimización bayesiana son las referencias de los procesos reales. La metodología sigue las guías de Montgomery~\cite[§1.5]{montgomery} y Garnett~\cite[cap.~9]{garnett}.

\section{Del problema al optimizador}

\begin{algoritmo}[Lista de diseño de un optimizador experimental]\label{alg:lista}
\begin{enumerate}[label=\arabic*.]
  \item \emph{Objetivo}: una cifra con unidades y su justificación económica (coste por pieza buena, emisiones por m$^3$). Si hay varias, decidir si se escalariza o se busca el frente de Pareto.
  \item \emph{Variables}: las que se pueden fijar en planta, con rangos seguros acordados con quien opera el proceso; escala lineal o logarítmica.
  \item \emph{Restricciones}: clasificarlas (conocidas, caras, ocultas, de seguridad) según el \cref{tema:1}.
  \item \emph{Ruido}: de dónde viene, cuánto vale (réplicas o un estudio previo) y si es heterocedástico.
  \item \emph{Presupuesto y logística}: número de ensayos, cuántos en paralelo, cuánto tarda cada uno y si hay fidelidades más baratas.
  \item \emph{Cribado}: un factorial o una fracción para descartar variables y detectar interacciones.
  \item \emph{Bucle}: diseño inicial, sustituto, adquisición, recomendación (\cref{alg:bo}).
  \item \emph{Parada y confirmación}: criterio de parada (\cref{tema:6}) y ensayos repetidos en la recomendación.
  \item \emph{Registro}: guardar cada ensayo con su contexto (lote de material, turno, operario) para poder reutilizarlo.
\end{enumerate}
\end{algoritmo}

\addtocounter{figure}{1}%

La comparación con la práctica tradicional de variar un factor cada vez (OFAT) muestra por qué se necesita todo lo anterior.

\addtocounter{definicion}{1}%

\section{Caso A: el ajuste completo de la inyectora}

\addtocounter{definicion}{1}%

\section{Caso B: un hormigón con menos CO\texorpdfstring{$_2$}{2}}

El cemento es responsable de una fracción importante de las emisiones globales de CO$_2$, lo que motiva sustituir parte de él por adiciones como las cenizas volantes~\cite{ament}. La dificultad es que la resistencia, que se mide rompiendo probetas a los 28 días, baja con la sustitución.

\begin{definicion}[Resistencia característica y resistencia media objetivo]\label{def:fck}
La resistencia característica $f_{ck}$ es el cuantil del $5\,\%$ de la resistencia a compresión a 28 días. Si la resistencia de las amasadas de una mezcla es normal con media $f_{cm}$ y desviación $\sigma_p$ (variabilidad de producción), $f_{ck}=f_{cm}-1{,}645\,\sigma_p$, y para garantizar $f_{ck}$ la mezcla debe tener $f_{cm}\ge f_{ck}+1{,}645\,\sigma_p$.
\end{definicion}

La fórmula es el cuantil de la normal (módulo~0.8, §6.3): $\Pr[f<f_{cm}-1{,}645\sigma_p]=\Phi(-1{,}645)=0{,}05$.

\addtocounter{definicion}{1}%

\begin{observacion}[Objetivo conocido, restricción desconocida]\label{obs:objconocido}
Cuando el objetivo $g$ es una fórmula conocida y lo caro es la restricción, la mejora en un candidato es determinista, $I(x)=\max(g^*-g(x),0)$, y la EI restringida de la \cref{prop:eic} se reduce a
\[
\alpha(x)=\max\bigl(g^*-g(x),0\bigr)\cdot\Pr\bigl[f_{28}(x)\ge 36{,}58\bigr]=\max\bigl(g^*-g(x),0\bigr)\cdot\Phi\Bigl(\frac{\mu_{28}(x)-36{,}58}{\sigma_{28}(x)}\Bigr).
\]
Solo hay que modelar lo que no se conoce. La misma idea mejora el caso de la inyectora, cuyo tiempo de ciclo es conocido: basta modelar la fracción defectuosa.
\end{observacion}

\addtocounter{definicion}{1}%

\addtocounter{definicion}{1}%

\begin{table}[ht]
\centering\small
\begin{tabularx}{\textwidth}{@{}>{\raggedright\arraybackslash}p{2.8cm}>{\raggedright\arraybackslash}X>{\raggedright\arraybackslash}p{2.2cm}>{\raggedright\arraybackslash}p{3.3cm}@{}}
\toprule
Herramienta & Para qué sirve & Licencia o tipo & Cuándo elegirla \\
\midrule
Ax, interfaz \texttt{Client}~\cite{ax} & Pedir ensayos (\texttt{get\_next\_trials}), registrar resultados (\texttt{complete\_trial}), restricciones de resultado, ensayos pendientes & MIT & Campañas en planta o laboratorio con varias personas \\
BoTorch~\cite{botorch} & Adquisiciones con objetivos conocidos, restricciones y multifidelidad & MIT & Formulaciones a medida como la \cref{obs:objconocido} \\
pyDOE3, \texttt{scipy.stats.qmc} & Cribado y diseño inicial & BSD & Fases 6 y 7 de la lista \\
Sustainable\-Concrete~\cite{ament} & Código de los modelos y del bucle de Ament \emph{et al.} & Código abierto (GitHub) & Punto de partida para mezclas reales \\
Hoja de cálculo o base de datos & Registro de ensayos con contexto & — & Siempre: es el activo más valioso \\
\bottomrule
\end{tabularx}
\caption{Tecnologías para desplegar un optimizador experimental.}
\label{tab:tec5}
\end{table}

\addtocounter{definicion}{1}%

